\begin{frame}{역사: ``Temporal-Difference'' 이름의 유래}
  \begin{itemize}
    \item 이름은 갱신식 구조를 그대로 번역: 추정치가 \textbf{시간 방향}으로
      \textbf{차이}를 만드는 것 --- $V(s')$(시각 $t{+}1$)과 $V(s)$(시각 $t$)의
      \textbf{차}가 갱신을 이끈다.
    \item MC = ``현재 추정치'' vs ``에피소드 끝의 실제 리턴'',
      DP = ``현재 추정치'' vs ``정확한 기댓값''.
    \item TD는 ``현재 추정치'' vs \textbf{한 스텝 앞의 추정치} ---
      ``시간의 차이''가 목표값 재료가 되는 \textbf{유일한} 방법.
  \end{itemize}
  \vskip0.6em
  \begin{block}{1983, mountain car}
    Barto·Sutton·Anderson(1983)의 *Neural learning of the optimal way to
    climb a mountain car*에서 처음 제안.
    \textbf{mountain car} 문제: 계곡 바닥의 자동차가 모멘텀을 만들어 언덕
    중 어느 쪽이든 올라가게. 물리 모델이 복잡해 DP로 정확히 풀기 어렵고,
    MC처럼 에피소드를 기다리면 너무 느린 --- \textbf{정확히 TD가 필요한}
    환경.
  \end{block}
  \vskip0.4em
  1984년 Sutton 박사논문(CMU, 지도 A. Barto)이 체계화 $\to$ 1989년
  Watkins의 Q-learning $\to$ 2013년 DQN(Week 9) --- 모든 off-policy TD
  학습의 뼈대가 이 ``한 스텝의 차이''에 놓여 있다.
\end{frame}
